% Time-transfer measurement over the whole corpus, split by recorder build.
% Every number here is produced by scripts/tools/sync_analysis.py from the released
% sessions; nothing is hand-entered. Regenerate with:
%     python3 scripts/tools/sync_analysis.py --csv sync_per_session.csv

\subsection{Time-Transfer Accuracy Over the Whole Corpus}\label{sec:sync_measured}

The purpose of the shared timebase is to answer one question: given camera frame $N$,
which inertial sample is the same physical instant? This section measures how well that
question can be answered, on every session rather than a subset.

\paragraph{What is paired.} The camera drives its frame-sync pulse into the FSYNC pad of
the bar-mounted IMU, which marks the coincident sample. A tagged sample is therefore a
camera exposure expressed in the \emph{bar's} clock, and the camera's own frame timestamp
is the same exposure in the \emph{camera's} clock: one physical instant, two clocks. Pairs
were re-derived offline from the stored streams. Pulses tagged twice within half a frame
period are merged, and the pairing walk is initialised from the ratio of the two total
spans rather than from a leading window --- the two oscillators differ by hundreds of
parts per million, so over a minute they separate by several frame periods and a slope
fitted to the first seconds mispredicts the end of the session by more than the matching
gate.

\paragraph{The recorder changed once.} The corpus was collected on four days with two
recorder builds, and the session metadata records which. Build \texttt{d225757} was used
on 10 and 11 May and for part of 17 May; \texttt{cdecba0} for the remainder of 17 May, 18
May and 20 May. Rather than exclude the earlier sessions, Table~\ref{tab:sync_by_build}
reports the two groups separately.

\begin{table}[t]
  \centering
  \caption{Time transfer, by recorder build. All 84 sessions, 382\,102 pulse--frame pairs.
           The residual is also given in inertial samples, since that is the quantity a
           consumer of the corpus needs: how far the correspondence between a camera frame
           and an inertial sample can be wrong.}
  \label{tab:sync_by_build}
  \begin{tabular}{lcc}
    \hline
    & \texttt{d225757} & \texttt{cdecba0} \\
    & 10--17 May & 17--20 May \\
    \hline
    Sessions                        & 48      & 36      \\
    Pulse--frame pairs              & 193\,794 & 188\,308 \\
    \hline
    Median $|$error$|$ (ms)         & 1.43    & 1.13    \\
    95th percentile (ms)            & 4.64    & 4.00    \\
    99th percentile (ms)            & 5.37    & 5.20    \\
    Worst case (ms)                 & 6.85    & 6.11    \\
    \hline
    Median (inertial samples)       & 1.4     & 1.1     \\
    99th percentile (samples)       & 5.4     & 5.2     \\
    \hline
    \multicolumn{3}{l}{\emph{Oscillator rate error removed by the anchor}} \\
    Median $|a-1|$ (ppm)            & 337     & 74      \\
    90th percentile (ppm)           & 944     & 414     \\
    \hline
  \end{tabular}
\end{table}

The later build is better by roughly a third of an inertial sample and by a factor of four
in the rate error it has to remove. Neither difference changes what the corpus supports:
in both groups a camera frame maps to an inertial sample to within about one sample
typically and five samples at the 99th percentile. At the fastest point of a lift the bar
travels \SI{1.8}{\milli\metre} per millisecond, so a one-sample misalignment contributes
about \SI{5}{\milli\metre\per\second} of apparent velocity disagreement --- twenty-five
times smaller than the \SI{127}{\milli\metre\per\second} RMSE of the inertial baseline the
corpus exists to evaluate. The synchronisation is therefore not the limiting term in any
agreement figure reported here, and that is now measured rather than assumed.

\subsection{Which Clock Carries the Residual}\label{sec:sync_budget}

The residual can be attributed by comparing each stream against a uniform
\SI{90}{\hertz} cadence on its own terms. This must be indexed by the camera's hardware
frame counter and not by row: a dropped frame is absent from the file, so indexing by row
makes every subsequent frame appear one period late and inflates the figure roughly
tenfold.

\begin{table}[t]
  \centering
  \caption{Deviation from a uniform \SI{90}{\hertz} cadence, RMS in the median session.
           The pulse train is what the camera emitted, read in the bar's clock; the
           timestamps are what the camera reports about those same exposures.}
  \label{tab:cadence}
  \begin{tabular}{lcc}
    \hline
    & \texttt{d225757} & \texttt{cdecba0} \\
    \hline
    Pulse train, in the bar's clock (ms) & 0.293 & 0.293 \\
    Camera's own timestamps (ms)         & 1.372 & 1.353 \\
    \hline
  \end{tabular}
\end{table}

The pulse-train figure is the inertial sampling quantisation and nothing else: a tag is
resolved to the \SI{1.012}{\milli\second} sampling grid, contributing
$1.012/\sqrt{12}=\SI{292}{\micro\second}$, and \SI{293}{\micro\second} is measured. It is
the same in every session of the corpus (range 0.291--0.293), across both builds and all
four days. The camera's timestamps are \emph{4.6 times noisier than the pulse train that
produced the very frames they label}, and in 7 of 84 sessions the camera's clock
misbehaves outright, reaching \SI{48}{\milli\second} RMS and \SI{5524}{ppm} of period
error against the inertial train's steady \SI{11.1278}{\milli\second}.

This has a practical consequence beyond the error budget. Because every exposure is
already tagged inside the inertial stream, the frame times are better taken from the IMU
than from the camera. The hardware trigger is not only a means of measuring alignment; it
is a means of replacing the noisier of the two clocks.

\subsection{Dropped Frames and the Reference Timebase}\label{sec:dropped_frames}

The camera does not deliver every frame it exposes. Across the corpus 838 frames are
missing, a median of 7 per session and 82 in the worst, identified by discontinuities in
the camera's hardware frame counter. A dropped frame is simply absent from the marker
file, so any pipeline that treats the row index as a timebase closes the hole up: it
believes \SI{11.1}{\milli\second} elapsed where \SI{22.2}{\milli\second} did, and every
subsequent event is early --- by \SI{0.91}{\second} at the end of the worst session.

Because the pulse train records an exposure for frames the camera then failed to deliver,
the drops are locatable to the sample and the reference track can be reconstructed on the
camera's true frame grid, with a missing frame treated exactly as a frame on which the
marker was not seen: an instant with no measurement, where the smoother must supply the
motion and report its own uncertainty for doing so.

Correcting this leaves the annotation structurally unchanged --- the same 1409
repetitions, no session changing its count, and 92\% of boundaries landing on the same
physical frame --- while changing the kinematics where it matters. The median correction
is zero, but the worst peak concentric velocity moves by
\SI{161}{\milli\metre\per\second}, the worst mean concentric velocity by
\SI{63}{\milli\metre\per\second}, and the worst repetition duration by
\SI{322}{\milli\second}. Every correction is downward in velocity, since a compressed
timebase inflates it. The largest is greater than the RMSE of the inertial baseline the
reference is used to judge, so it is not a rounding matter, and the number of repetitions
flagged as containing an unmeasured instant rises from 14 to 369 --- information a
consumer of the corpus needs and did not previously have.
